% ch3_a 曲率 (书页 93-105 / p108-p120)
% 第3章 Curvature
\chapter{曲率}

\section{概述}

我们都知道曲率是什么意思，至少是非正式地知道；而且在本书前两章里，我们一直随意地不时提及曲率这个概念，却从未认真定义过它。显然，曲率以某种方式依赖于度规（metric），度规定义了我们流形的几何；但如何把曲率归属于任何一个给定的度规，并不是一目了然的（因为正如我们所见，即使是平直空间的度规，在一个足够“铺张”的坐标系里也可以显得任意复杂）。数学中常见的情况是：要把我们关于一个概念的直觉形式化为可用的数学结构，需要相当多的细心；把我们所想的“曲率”形式化，正是本章的主题。

我们即将发展出来的技巧对本学科至关重要；可以放心地说，这一章每页有用公式的密度比其他任何一章都高。让我们快速总结一下其中最重要的那些，为接下来的形式体系提供一张路线图。

曲率表现自身的所有方式都依赖于一种叫做“联络”（connection）的东西，它为我们提供了一种联系相邻点的切空间中矢量的途径。存在一个唯一的、可以由度规构造出来的联络，它浓缩在一个叫做 Christoffel 符号（Christoffel symbol）的对象之中：
\begin{equation*}
\Gamma^\lambda_{\mu\nu} = \frac{1}{2} g^{\lambda\sigma}\left(\partial_\mu g_{\nu\sigma} + \partial_\nu g_{\sigma\mu} - \partial_\sigma g_{\mu\nu}\right).
\tag{3.1}
\end{equation*}

这种记号让 $\Gamma^\lambda_{\mu\nu}$ 看起来像一个张量，但事实上它不是；这就是为什么我们把它叫做一个“对象”（object）或“符号”（symbol）。联络最基本的用途是用来取协变导数（covariant derivative）$\nabla_\mu$（偏导数的推广）；矢量场 $V^\nu$ 的协变导数为
\begin{equation*}
\nabla_\mu V^\nu = \partial_\mu V^\nu + \Gamma^\nu_{\mu\sigma} V^\sigma\,,
\tag{3.2}
\end{equation*}
而其他种类的张量的协变导数也由类似的表达式给出。联络还出现在测地线（geodesic，直线概念的推广）的定义中。参数化曲线 $x^\mu(\lambda)$ 是测地线，如果它满足
\begin{equation*}
\frac{d^2 x^\mu}{d\lambda^2} + \Gamma^\mu_{\rho\sigma}\frac{d x^\rho}{d\lambda}\frac{d x^\sigma}{d\lambda} = 0,
\tag{3.3}
\end{equation*}
这就是著名的测地线方程（geodesic equation）。

最后，曲率的技术性表述包含在 Riemann 曲率张量（Riemann tensor）之中，它是一个 $(1,3)$ 型张量，由联络构造出来：
\begin{equation*}
R^\rho{}_{\sigma\mu\nu} = \partial_\mu \Gamma^\rho_{\nu\sigma} - \partial_\nu \Gamma^\rho_{\mu\sigma} + \Gamma^\rho_{\mu\lambda}\Gamma^\lambda_{\nu\sigma} - \Gamma^\rho_{\nu\lambda}\Gamma^\lambda_{\mu\sigma}.
\tag{3.4}
\end{equation*}
我们想了解的关于流形曲率的一切都由 Riemann 张量给出；当且仅当度规完全平直时它才为零。广义相对论的 Einstein 方程把这个张量的某些分量与能量-动量张量联系起来。

这四个公式在弯曲流形的研究中都具有头等重要性。接下来我们将会看到，它们是如何从对平直空间中那些熟悉的几何观念如何适应于这个更一般背景的细致考察中产生出来的。

\section{协变导数}

在我们对流形的讨论中，有一点已经很清楚：流形一旦定义，就有各种各样可以谈论的概念了——我们可以定义函数、求导数、考虑参数化路径、搭建张量，等等。而另一些概念，比如一个区域的体积或一条路径的长度，则需要额外的某种结构，也就是度规的引入。把曲率这个概念想成是只依赖于度规的东西，似乎是很自然的。然而，在更仔细的处理中我们会发现，曲率依赖于联络，而联络依赖于或不依赖于度规皆有可能。尽管如此，我们也将证明，度规的存在蕴含着一个唯一确定的联络，这个联络的曲率可以看作是该度规的曲率。这正是广义相对论中使用的那个联络，所以在这个特定背景下，把曲率看作刻画度规的性质而不引入任何附加结构，是合理的。

当我们试图处理“偏导数不是一个好的张量算符”这一问题时，联络就成为必要的了。我们想要的是一个协变导数（covariant derivative），即这样一个算符：在平直空间的惯性坐标下它还原为偏导数，而在任意流形上它按张量方式变换。按惯例，人们会花不少篇幅来论证引入协变导数的动机，但实际上这个需求是显而易见的：诸如 $\partial_\mu T^{\mu\nu} = 0$ 这样的方程总得设法推广到弯曲空间去。所以我们就一致同意：协变导数是个好东西，值得拥有，然后着手把它建立起来。

在惯性坐标下的平直空间中，偏导数算符 $\partial_\mu$ 是一个从 $(k,l)$ 型张量场到 $(k,l+1)$ 型张量场的映射，它对其自变量的作用是线性的，并且在张量积上满足 Leibniz 律。在我们现在想要考虑的更一般的情形下，这一切仍然成立，但偏导数所提供的这个映射依赖于所用的坐标系。因此我们想定义一个协变导数（covariant derivative）算符 $\nabla$ 来行使偏导数的职能，但要以一种与坐标无关的方式进行。与其直接假定答案（那样做其实完全可以接受），不如仔细想想偏导数的协变推广\emph{应该}具有些什么性质——数学结构毕竟是人类的发明，而不是躺在人行道上等着被捡到的。我们一开始先要求：$\nabla$ 是一个从 $(k,l)$ 型张量场到 $(k,l+1)$ 型张量场的映射，并且具有下面两条性质：
\begin{enumerate}
\item 线性性：$\nabla(T + S) = \nabla T + \nabla S$；
\item Leibniz（乘积）律：$\nabla(T \otimes S) = (\nabla T)\otimes S + T\otimes(\nabla S)$.
\end{enumerate}

如果 $\nabla$ 要服从 Leibniz 律，那么它总可以写成偏导数再加上某个线性变换的形式。也就是说，取协变导数时我们先取偏导数，然后施加一个修正，使结果成为协变的。【这个听起来挺合理的论断我们不打算证明；感兴趣的读者可参阅 Wald (1984)。】让我们考虑这对矢量 $V^\nu$ 的协变导数意味着什么。它意味着：对每个方向 $\mu$，协变导数 $\nabla_\mu$ 将由偏导数 $\partial_\mu$ 加上一个修正给出，该修正由一组 $n$ 个矩阵 $(\Gamma_\mu)^\rho{}_\sigma$ 规定（对每个 $\mu$ 有一个 $n\times n$ 矩阵，其中 $n$ 是流形的维数）。实际上这对括号通常被省掉，这些矩阵——即所谓的\emph{联络系数}（connection coefficients）——以一种随意的指标位置写成为 $\Gamma^\rho_{\mu\sigma}$。于是我们有
\begin{equation*}
\boxed{\;\nabla_\mu V^\nu = \partial_\mu V^\nu + \Gamma^\nu_{\mu\lambda} V^\lambda.\;}
\tag{3.5}
\end{equation*}

注意在第二项里，原来在 $V$ 上的指标跑到了 $\Gamma$ 上，同时有一个新的指标被求和。如果这就是用偏导数表出的矢量协变导数，那么我们应该能通过要求左边是一个 $(1,1)$ 型张量，来确定 $\Gamma^\nu_{\mu\lambda}$ 的变换性质。也就是说，我们希望变换规律是
\begin{equation*}
\nabla_{\mu'} V^{\nu'} = \frac{\partial x^\mu}{\partial x^{\mu'}}\frac{\partial x^{\nu'}}{\partial x^\nu}\nabla_\mu V^\nu.
\tag{3.6}
\end{equation*}
让我们先看左边；可以用式 (3.5) 把它展开，再变换那些我们已经理解的部分（也就是除了 $\Gamma^{\nu'}_{\mu'\lambda'}$ 之外的全部）：
\begin{align*}
\nabla_{\mu'} V^{\nu'} &= \partial_{\mu'} V^{\nu'} + \Gamma^{\nu'}_{\mu'\lambda'} V^{\lambda'}\\
&= \frac{\partial x^\mu}{\partial x^{\mu'}}\frac{\partial x^{\nu'}}{\partial x^\nu}\partial_\mu V^\nu + \frac{\partial x^\mu}{\partial x^{\mu'}}V^\nu\frac{\partial}{\partial x^\mu}\frac{\partial x^{\nu'}}{\partial x^\nu} + \Gamma^{\nu'}_{\mu'\lambda'}\frac{\partial x^{\lambda'}}{\partial x^\lambda}V^\lambda.
\tag{3.7}
\end{align*}
在右边，我们同样可以把 $\nabla_\mu V^\nu$ 展开：
\begin{equation*}
\frac{\partial x^\mu}{\partial x^{\mu'}}\frac{\partial x^{\nu'}}{\partial x^\nu}\nabla_\mu V^\nu = \frac{\partial x^\mu}{\partial x^{\mu'}}\frac{\partial x^{\nu'}}{\partial x^\nu}\partial_\mu V^\nu + \frac{\partial x^\mu}{\partial x^{\mu'}}\frac{\partial x^{\nu'}}{\partial x^\nu}\Gamma^\nu_{\mu\lambda}V^\lambda.
\tag{3.8}
\end{equation*}
最后这两个表达式应当相等；两边的第一项完全相同因而抵消，于是我们得到
\begin{equation*}
\Gamma^{\nu'}_{\mu'\lambda'}\frac{\partial x^\lambda}{\partial x^{\lambda'}}V^\lambda + \frac{\partial x^\mu}{\partial x^{\mu'}}V^\lambda\frac{\partial}{\partial x^\mu}\frac{\partial x^{\nu'}}{\partial x^\lambda} = \frac{\partial x^\mu}{\partial x^{\mu'}}\frac{\partial x^{\nu'}}{\partial x^\nu}\Gamma^\nu_{\mu\lambda}V^\lambda,
\tag{3.9}
\end{equation*}
这里我们把一个哑指标从 $\nu$ 改成了 $\lambda$。这个方程对任意矢量 $V^\lambda$ 都必须成立，所以我们可以在两边把它消去。然后，通过乘上 $\partial x^\lambda/\partial x^{\sigma'}$ 并把 $\sigma'$ 重新标记为 $\lambda'$，就可以把带撇坐标中的联络系数孤立出来。结果是
\begin{equation*}
\Gamma^{\nu'}_{\mu'\lambda'} = \frac{\partial x^\mu}{\partial x^{\mu'}}\frac{\partial x^\lambda}{\partial x^{\lambda'}}\frac{\partial x^{\nu'}}{\partial x^\nu}\Gamma^\nu_{\mu\lambda} + \frac{\partial x^\mu}{\partial x^{\mu'}}\frac{\partial x^\lambda}{\partial x^{\lambda'}}\frac{\partial^2 x^{\nu'}}{\partial x^\mu\partial x^\lambda}.
\tag{3.10}
\end{equation*}
当然，这并不是张量的变换规律——右边的第二项把它破坏了。不过没关系，因为\emph{联络系数并不是某个张量的分量}。它们是被有意构造成非张量的，但构造的方式使得式 (3.5) 这个组合按张量变换——偏导数变换中多出来的项与 $\Gamma$ 变换中多出来的项恰好抵消。正因如此，我们对联络系数的指标位置并不那么讲究；它们不是张量，所以你应当尽量不要对它们作指标的升降。

那么其他种类的张量的协变导数呢？通过与矢量情形类似的推理，1-形式（one-form）的协变导数同样可以表示为偏导数加上某个线性变换。但迄今为止还没有任何理由认为，表示这个变换的矩阵应当与系数 $\Gamma^\nu_{\mu\lambda}$ 有什么关系。一般地，我们可以写出类似
\begin{equation*}
\nabla_\mu\omega_\nu = \partial_\mu\omega_\nu + \widetilde{\Gamma}^\lambda_{\mu\nu}\,\omega_\lambda,
\tag{3.11}
\end{equation*}
其中 $\widetilde{\Gamma}^\lambda_{\mu\nu}$ 是对每个 $\mu$ 的一组新矩阵。请注意各种不同的指标各自都在什么位置。可以直接推导出，$\widetilde{\Gamma}$ 的变换性质必须与 $\Gamma$ 相类似，否则 $\nabla_\mu\omega_\nu$ 就不会按张量变换；但除此之外，两者之间没有任何已确立的关系。为了确立这种关系，除了上面那两条性质之外，我们还需要再引入两条我们希望协变导数具有的性质：
\begin{enumerate}
\setcounter{enumi}{2}
\item 与缩并可交换：$\nabla_\mu(T^\lambda{}_{\lambda\rho}) = (\nabla T)_\mu{}^\lambda{}_{\lambda\rho}$，
\item 在标量上还原为偏导数：$\nabla_\mu\phi = \partial_\mu\phi$.
\end{enumerate}

没有办法“推导”出这些性质；我们只是要求它们作为协变导数定义的一部分而成立。注意，性质 3 等价于说 Kronecker delta（恒等映射）是协变常数的（covariantly constant）：$\nabla_\mu\delta^\lambda_\sigma = 0$；这当然是一个合理的要求。

让我们看看这些新性质蕴含着什么。给定某个 1-形式场 $\omega_\mu$ 和矢量场 $V^\mu$，我们可以对由 $\omega_\lambda V^\lambda$ 定义的标量取协变导数，得到
\begin{align*}
\nabla_\mu(\omega_\lambda V^\lambda) &= (\nabla_\mu\omega_\lambda)V^\lambda + \omega_\lambda(\nabla_\mu V^\lambda)\\
&= (\partial_\mu\omega_\lambda)V^\lambda + \widetilde{\Gamma}^\sigma_{\mu\lambda}\omega_\sigma V^\lambda + \omega_\lambda(\partial_\mu V^\lambda) + \omega_\lambda\Gamma^\lambda_{\mu\rho}V^\rho.
\tag{3.12}
\end{align*}
但由于 $\omega_\lambda V^\lambda$ 是一个标量，它也必须可以由偏导数给出：
\begin{align*}
\nabla_\mu(\omega_\lambda V^\lambda) &= \partial_\mu(\omega_\lambda V^\lambda)\\
&= (\partial_\mu\omega_\lambda)V^\lambda + \omega_\lambda(\partial_\mu V^\lambda).
\tag{3.13}
\end{align*}
只有当式 (3.12) 中带联络系数的那些项彼此抵消时，这才可能成立；也就是说，重新安排哑指标后，我们必须有
\begin{equation*}
0 = \widetilde{\Gamma}^\sigma_{\mu\lambda}\omega_\sigma V^\lambda + \Gamma^\sigma_{\mu\lambda}\omega_\sigma V^\lambda.
\tag{3.14}
\end{equation*}
但 $\omega_\sigma$ 和 $V^\lambda$ 都是完全任意的，所以
\begin{equation*}
\widetilde{\Gamma}^\sigma_{\mu\lambda} = -\Gamma^\sigma_{\mu\lambda}.
\tag{3.15}
\end{equation*}
我们所施加的这两个附加条件，于是使我们能够用与矢量情形相同的联络系数来表达 1-形式的协变导数，只是现在带一个负号（并且指标的搭配方式略有不同）：
\begin{equation*}
\boxed{\;\nabla_\mu\omega_\nu = \partial_\mu\omega_\nu - \Gamma^\lambda_{\mu\nu}\,\omega_\lambda.\;}
\tag{3.16}
\end{equation*}

联络系数编码了对任意阶张量取协变导数所需的全部信息，这应当不足为怪。这个公式相当直截了当：对每个上指标，引入一个带单个 $+\Gamma$ 的项；对每个下指标，引入一个带单个 $-\Gamma$ 的项：
\begin{align*}
\nabla_\sigma T^{\mu_1\mu_2\cdots\mu_k}{}_{\nu_1\nu_2\cdots\nu_l} ={}& \partial_\sigma T^{\mu_1\mu_2\cdots\mu_k}{}_{\nu_1\nu_2\cdots\nu_l}\\
&+ \Gamma^{\mu_1}_{\sigma\lambda}T^{\lambda\mu_2\cdots\mu_k}{}_{\nu_1\nu_2\cdots\nu_l} + \Gamma^{\mu_2}_{\sigma\lambda}T^{\mu_1\lambda\cdots\mu_k}{}_{\nu_1\nu_2\cdots\nu_l} + \cdots\\
&- \Gamma^{\lambda}_{\sigma\nu_1}T^{\mu_1\mu_2\cdots\mu_k}{}_{\lambda\nu_2\cdots\nu_l} - \Gamma^{\lambda}_{\sigma\nu_2}T^{\mu_1\mu_2\cdots\mu_k}{}_{\nu_1\lambda\cdots\nu_l} - \cdots.
\tag{3.17}
\end{align*}

这就是协变导数的一般表达式。你可以自行验证它；它来自我们所建立的这组公理，再加上“各类张量都是与坐标无关的实体”这一通常要求。有时人们使用另一种记号：正如用逗号表示偏导数，用分号表示协变导数：
\begin{equation*}
\nabla_\sigma T^{\mu_1\mu_2\cdots\mu_k}{}_{\nu_1\nu_2\cdots\nu_l} \equiv T^{\mu_1\mu_2\cdots\mu_k}{}_{\nu_1\nu_2\cdots\nu_l;\sigma}.
\tag{3.18}
\end{equation*}
再一次强调，在本书中我们将坚持使用“$\nabla_\sigma$”。
那么，要定义一个协变导数，我们需要在流形上放一个联络，它在某个坐标系中由一组系数 $\Gamma^\lambda_{\mu\nu}$ 规定（在 $n=4$ 维时有 $n^3=64$ 个独立分量），这些系数按式 (3.10) 变换。（“联络”（connection）这个名字来源于这样一个事实：它被用来把矢量从一个切空间输运到另一个切空间，我们很快就会看到；它有时用来指算符 $\nabla$，有时用来指系数 $\Gamma^\lambda_{\mu\nu}$。）显然，我们可以在任何流形上定义出许许多多联络，而每一个都蕴含着一种互不相同的协变微分的概念。在广义相对论中，这种自由度不是什么大问题，因为事实表明每个度规都定义一个唯一的联络，这正是 GR 中使用的那一个。让我们看看这是如何做到的。

首先要注意到，两个联络之差是一个张量。设想我们定义了两种不同的协变导数 $\nabla_\mu$ 和 $\widehat{\nabla}_\mu$，分别带有相联系的联络系数 $\Gamma^\lambda_{\mu\nu}$ 和 $\widehat{\Gamma}^\lambda_{\mu\nu}$。那么差
\begin{equation*}
S^\lambda{}_{\mu\nu} = \Gamma^\lambda_{\mu\nu} - \widehat{\Gamma}^\lambda_{\mu\nu}
\tag{3.19}
\end{equation*}
是一个 $(1,2)$ 型张量。（注意我们不得不为指标的摆放选定一个约定。）我们可以用蛮力来证明这一点，即代入联络系数的变换规律，但让我们稍微俏皮一点。给定一个任意矢量场 $V^\lambda$，我们知道 $\nabla_\mu V^\lambda$ 和 $\widehat{\nabla}_\mu V^\lambda$ 都是张量，因此它们的差也必是张量。这个差就是
\begin{align*}
\nabla_\mu V^\lambda - \widehat{\nabla}_\mu V^\lambda &= \partial_\mu V^\lambda + \Gamma^\lambda_{\mu\nu}V^\nu - \partial_\mu V^\lambda - \widehat{\Gamma}^\lambda_{\mu\nu}V^\nu\\
&= S^\lambda{}_{\mu\nu}V^\nu.
\tag{3.20}
\end{align*}
由于 $V^\lambda$ 是任意的，而左边是张量，所以 $S^\lambda{}_{\mu\nu}$ 必定是张量。作为一个平凡的推论，我们得知任何一组联络系数都可以表示为某个基准联络（fiducial connection）加上一个张量性的修正：
\begin{equation*}
\Gamma^\lambda_{\mu\nu} = \widehat{\Gamma}^\lambda_{\mu\nu} + S^\lambda{}_{\mu\nu}.
\tag{3.21}
\end{equation*}

接下来注意，给定一个由 $\Gamma^\lambda_{\mu\nu}$ 规定的联络，我们可以仅仅通过交换下指标立即构造出另一个联络。也就是说，系数集合 $\Gamma^\lambda_{\nu\mu}$ 也将按式 (3.10) 变换（因为出现在最后一项中的偏导数可以交换次序），所以它们确定一个互不相同的联络。这样一来，就有了一个可以与任何给定联络相联系的张量，即所谓的\emph{挠率张量}（torsion tensor），其定义为
\begin{equation*}
T^\lambda{}_{\mu\nu} = \Gamma^\lambda_{\mu\nu} - \Gamma^\lambda_{\nu\mu} = 2\Gamma^\lambda_{[\mu\nu]}.
\tag{3.22}
\end{equation*}
显然，挠率在其下指标上是反对称的，而在下指标上对称的联络被称为“无挠的”（torsion-free）。

现在我们可以在一个带有度规 $g_{\mu\nu}$ 的流形上，通过引入两条附加性质来定义一个唯一的联络：
\begin{itemize}
\item 无挠：$\Gamma^\lambda_{\mu\nu} = \Gamma^\lambda_{(\mu\nu)}$；
\item 度规相容性（metric compatibility）：$\nabla_\rho g_{\mu\nu} = 0$.
\end{itemize}

如果一个联络对度规取的协变导数处处为零，就称这个联络是\emph{度规相容}的（metric compatible）。这蕴含着几个不错的性质。第一，容易证明，Levi–Civita 张量和逆度规的协变导数同样为零：
\begin{align*}
\nabla_\lambda\epsilon_{\mu\nu\rho\sigma} &= 0\\
\nabla_\rho g^{\mu\nu} &= 0.
\tag{3.23}
\end{align*}
第二，度规相容的协变导数与指标的升降可以交换次序。于是，对某个矢量场 $V^\lambda$，
\begin{equation*}
g_{\mu\lambda}\nabla_\rho V^\lambda = \nabla_\rho(g_{\mu\lambda}V^\lambda) = \nabla_\rho V_\mu.
\tag{3.24}
\end{equation*}
如果联络不是度规相容的，那么在取协变导数时我们就得对指标的摆放非常小心。

因此，我们的论断是：在一个给定的流形上，与该流形上某个给定度规相容的无挠联络恰好只有一个。我们并不想把这两条要求纳入协变导数的定义；它们只不过是从许多可能的联络中挑出了一个。

我们可以通过推导出联络系数用度规表示的一个明显唯一的表达式，来同时证明存在性与唯一性。为此，我们把度规相容性方程按指标的三种不同排列展开：
\begin{align*}
\nabla_\rho g_{\mu\nu} &= \partial_\rho g_{\mu\nu} - \Gamma^\lambda_{\rho\mu}g_{\lambda\nu} - \Gamma^\lambda_{\rho\nu}g_{\mu\lambda} = 0\\
\nabla_\mu g_{\nu\rho} &= \partial_\mu g_{\nu\rho} - \Gamma^\lambda_{\mu\nu}g_{\lambda\rho} - \Gamma^\lambda_{\mu\rho}g_{\nu\lambda} = 0\\
\nabla_\nu g_{\rho\mu} &= \partial_\nu g_{\rho\mu} - \Gamma^\lambda_{\nu\rho}g_{\lambda\mu} - \Gamma^\lambda_{\nu\mu}g_{\rho\lambda} = 0.
\tag{3.25}
\end{align*}
从第一式中减去第二式和第三式，并利用联络的对称性，得到
\begin{equation*}
\partial_\rho g_{\mu\nu} - \partial_\mu g_{\nu\rho} - \partial_\nu g_{\rho\mu} + 2\Gamma^\lambda_{\mu\nu}g_{\lambda\rho} = 0.
\tag{3.26}
\end{equation*}
乘上 $g^{\rho\sigma}$ 之后即可直接解出联络。结果是
\begin{equation*}
\boxed{\;\Gamma^\sigma_{\mu\nu} = \frac{1}{2}g^{\sigma\rho}\left(\partial_\mu g_{\nu\rho} + \partial_\nu g_{\rho\mu} - \partial_\rho g_{\mu\nu}\right).\;}
\tag{3.27}
\end{equation*}
这个公式是本学科中最重要的公式之一；请把它牢牢记在心里。当然，我们只是证明了：如果度规相容且无挠的联络存在，它就必须具有式 (3.27) 的形式；你可以自行验证，式 (3.27) 的右边的确像联络一样变换。
我们从度规推导出的这个联络，正是传统的广义相对论所依据的联络。它有着不同的名字：有时叫 \emph{Christoffel 联络}，有时叫 \emph{Levi–Civita 联络}，有时叫 \emph{Riemann 联络}。与之相联系的联络系数有时被称为 \emph{Christoffel 符号}（Christoffel symbols），并写成 $\left\{^{\sigma}_{\mu\nu}\right\}$ 的样子；我们有时也会把它们叫做 Christoffel 符号，但不会使用这种古怪的记号。研究带有度规及其相应联络的流形的学问叫做黎曼几何（Riemannian geometry）；当度规具有洛伦兹号差时，有时也称为伪黎曼几何（pseudo-Riemannian geometry）。

在让我们的协变导数开工之前，应该先提几点杂七杂八的性质。首先，注意在通常的平直空间里，其实存在一个我们无时无刻不在使用的隐式联络——由平直度规构造出来的 Christoffel 联络。平直空间中 Christoffel 联络的系数在笛卡尔坐标下为零，但在曲线坐标系下则不然。以极坐标下的平面为例，其度规为
\begin{equation*}
ds^2 = dr^2 + r^2d\theta^2.
\tag{3.28}
\end{equation*}
逆度规的非零分量很容易求出，为 $g^{rr} = 1$ 和 $g^{\theta\theta} = r^{-2}$。注意，我们用 $r$ 和 $\theta$ 作指标，记号的含义一目了然。我们可以计算一个典型的联络系数：
\begin{align*}
\Gamma^r_{rr} &= \frac{1}{2}g^{r\rho}(\partial_r g_{r\rho} + \partial_r g_{\rho r} - \partial_\rho g_{rr})\\
&= \frac{1}{2}g^{rr}(\partial_r g_{rr} + \partial_r g_{rr} - \partial_r g_{rr})\\
&\quad + \frac{1}{2}g^{r\theta}(\partial_r g_{r\theta} + \partial_r g_{\theta r} - \partial_\theta g_{rr})\\
&= \frac{1}{2}(1)(0 + 0 - 0) + \frac{1}{2}(0)(0 + 0 - 0)\\
&= 0.
\tag{3.29}
\end{align*}
遗憾的是，它等于零。但并非所有系数都如此：
\begin{align*}
\Gamma^r_{\theta\theta} &= \frac{1}{2}g^{r\rho}(\partial_\theta g_{\theta\rho} + \partial_\theta g_{\rho\theta} - \partial_\rho g_{\theta\theta})\\
&= \frac{1}{2}g^{rr}(\partial_\theta g_{\theta r} + \partial_\theta g_{r\theta} - \partial_r g_{\theta\theta})\\
&= \frac{1}{2}(1)(0 + 0 - 2r)\\
&= -r.
\tag{3.30}
\end{align*}
继续埋头摇计算的手柄，最终我们发现
\begin{align*}
\Gamma^r_{\theta r} &= \Gamma^r_{r\theta} = 0\\
\Gamma^\theta_{rr} &= 0\\
\Gamma^\theta_{r\theta} &= \Gamma^\theta_{\theta r} = \frac{1}{r}\\
\Gamma^\theta_{\theta\theta} &= 0.
\tag{3.31}
\end{align*}
从这些以及类似的表达式出发，我们可以推导出曲线坐标系下发散、梯度和旋度的公式。

反过来说，即使在弯曲空间中，让 Christoffel 符号在某个单独的点上为零也仍然是可能的。这是因为，正如我们在上一章论证过的，总可以让度规的一阶导数在某点为零；于是由式 (3.27)，从这个度规导出的联络系数也将为零。当然，这只能在一点处成立，而不能在该点的某个邻域内都成立。我们将在 3.4 节更充分地讨论这一点。

另一个有用的性质是，矢量的散度（关于 Christoffel 联络）的公式有一个简化形式。$V^\mu$ 的协变散度为
\begin{equation*}
\nabla_\mu V^\mu = \partial_\mu V^\mu + \Gamma^\mu_{\mu\lambda}V^\lambda.
\tag{3.32}
\end{equation*}
容易证明 Christoffel 联络满足
\begin{equation*}
\Gamma^\mu_{\mu\lambda} = \frac{1}{\sqrt{|g|}}\partial_\lambda\sqrt{|g|},
\tag{3.33}
\end{equation*}
于是我们得到
\begin{equation*}
\nabla_\mu V^\mu = \frac{1}{\sqrt{|g|}}\partial_\mu\left(\sqrt{|g|}\,V^\mu\right).
\tag{3.34}
\end{equation*}
对更高阶张量的散度也有相应公式，但一般就没有这么大的简化了。

在 Stokes 定理的弯曲空间版本（见附录 E）中，我们使用的正是 Christoffel 协变导数。若 $V^\mu$ 是区域 $\Sigma$（其边界为 $\partial\Sigma$）上的矢量场，则 Stokes 定理为
\begin{equation*}
\boxed{\;\int_\Sigma \nabla_\mu V^\mu\sqrt{|g|}\,d^nx = \int_{\partial\Sigma} n_\mu V^\mu\sqrt{|\gamma|}\,d^{\,n-1}x,\;}
\tag{3.35}
\end{equation*}
其中 $n_\mu$ 垂直于 $\partial\Sigma$，而 $\gamma_{ij}$ 是 $\partial\Sigma$ 上的诱导度规。如果联络不是度规相容的或不是无挠的，这个方程里就会多出一些附加的项。

最后需要提到的一点是：为了构造定义良好的张量，把偏导数换成协变导数并不总是必要的；特别是，外微分和矢量场的对易子都可以仅用偏导数良好地定义，这本质上是因为两者都涉及一个反对称化操作，它恰好抵消了偏导数变换规律中的非张量部分。同一特征也意味着，反过来，它们同样可以很好地用（无挠的）协变导数来定义；反对称化使联络系数的项归于消失。于是，若 $\nabla$ 是 Christoffel 联络，$\omega_\mu$ 是一个 1-形式，而 $X^\mu$ 和 $Y^\mu$ 是矢量场，则我们
可以写成
\begin{equation*}
(d\omega)_{\mu\nu} = 2\partial_{[\mu}\omega_{\nu]} = 2\nabla_{[\mu}\omega_{\nu]}
\tag{3.36}
\end{equation*}
以及
\begin{equation*}
[X,Y]^\mu = X^\lambda\partial_\lambda Y^\mu - Y^\lambda\partial_\lambda X^\mu = X^\lambda\nabla_\lambda Y^\mu - Y^\lambda\nabla_\lambda X^\mu.
\tag{3.37}
\end{equation*}
如果联络不是无挠的，这些表达式中的最后一个等号就不再成立；外微分和对易子更基本的定义是那些用偏导数表述的定义。

在继续往下讲之前，让我们回顾一下我们一步步给数学构造添加结构的过程。我们从“集合”这个基本概念出发，假定你对此已经熟悉（至少非正式地熟悉，如果不严格的话）。我们引入了集合的开子集的概念；这等价于引入一个拓扑，从而把集合提升为拓扑空间。然后，通过要求每个开集看起来都像 $\mathbf{R}^n$ 的一个区域（对每个集合 $n$ 相同），并且要求坐标卡片彼此光滑地缝合在一起，拓扑空间就变成了流形。流形同时是一种非常灵活又非常强有力的结构，它自然而然地配备了切丛、各种阶的张量丛、取外微分的能力，等等。接着我们继续在流形上放置度规，得到带度规的流形（有时称为黎曼流形）。我们发现，独立于度规也可以引入联络，从而使我们能够取协变导数。然而，一旦有了度规，就自动存在一个唯一的无挠且度规相容的联络。同样，我们也可以引入一个独立的体积形式，尽管它其实由度规自动确定。原则上，没有什么能阻止我们在任何一个给定流形上引入不止一个联络、或不止一个体积形式、或不止一个度规。在广义相对论中，我们确实拥有一个物理的度规，它同时决定体积和协变导数，所以这些概念的相互独立并不是一个关键特征。

\section{平行移动与测地线}

既然我们已经知道如何取协变导数，就让我们退后一步，把它放进更一般的微分框架中来考察。我们把导数看作量化“某个东西变化得有多快”的手段。对张量而言，关键问题是“相对于什么变化？”一个普通函数在时空每一点定义一个数，而比较两个不同的数是直截了当的，所以我们不必惊讶函数的偏导数在任意流形上依然有效。但张量是从矢量和 对偶矢量到实数的映射，要比较时空不同点上的这种映射，就不清楚该怎么做了。既然我们已经成功地构造了协变导数，能否把它看作某种意义上在量度张量的变化率？答案证明是肯定的：协变导数量度的正是张量场的瞬时变化率——与该张量倘若被“平行移动”将会成为的样子相比较。

%FIG{3.1}: {在平直空间中，只需让矢量的笛卡尔（Cartesian）分量保持不变，我们就可以对矢量作平行移动。}

换句话说，联络定义了一种让张量（沿某条路径）保持不变的具体方式，我们正是以此为基准来比较相邻的张量。

结果证明，平行移动这个概念本身就很有意思，值得花些时间思考。回忆一下，在平直空间中，我们不必太在意“矢量是定义在各个单独的点上的切空间的元素”这个事实；在平直空间中比较不同点上的矢量实际上是非常自然的（这里“比较”指的是相加、相减、取点积等等）。之所以自然，是因为在平直空间中，把一个矢量从一点移动到另一点同时保持它不变是有意义的，如图 3.1 所示。这样一来，一旦我们把矢量从一点搬到了另一点，就可以进行矢量空间中通常允许的那些运算。

把一个矢量沿一条路径移动而始终保持不变，这个概念就叫做\emph{平行移动}（parallel transport）。平行移动需要一个联络才能良好定义；在平直空间中对矢量的直观操作，暗含地使用了这个空间上的 Christoffel 联络。平直空间与弯曲空间的关键差别在于：在弯曲空间中，\emph{把一个矢量从一点平行移动到另一点，其结果将依赖于这两点之间所走的路径。}不必先把平行移动的完整机制拼装出来，我们就可以借助关于二维球面的直觉看清这一点。从赤道上的一个矢量开始，让它指向一条恒定经线的方向。用显而易见的方式沿一条经线把它平行移动到北极。然后取原来的矢量，先沿赤道把它平行移动一个角度 $\theta$，再像刚才那样把它移到北极。如图 3.2 所示，显然这个矢量沿两条路径平行移动后，到达了同一个目的地却带着两个不同的值（相差一个转动角 $\theta$）。

这样看来，似乎没有什么自然的方法能把一个矢量唯一地从一个切空间移动到另一个切空间；我们总可以对它作平行移动，但结果依赖于路径，而且走哪条路径也没有自然的选择。
%FIG{3.2}: {二维球面上的平行移动。在弯曲流形上，平行移动的结果可能依赖于所走的路径。}

与我们遇到过的某些问题不同，\emph{这个问题无解}——我们干脆必须学会与如下事实共处：只有当两个矢量属于同一切空间时，才能以自然的方式比较它们。例如，互相擦肩而过的两个粒子有良好定义的相对速度，它不能超过光速。但在弯曲流形上不同两点处的两个粒子却没有任何良好定义的相对速度概念——这个概念根本就没有意义。当然，在某些特殊情形下，把它说得好像有意义仍然是有用的，但偶然的用处不能代替严格的定义。例如在宇宙学中，来自遥远星系的光相对于我们从近处静止光源将会观测到的频率发生了红移。由于这个现象与相对运动引起的传统多普勒效应极为相似，我们非常想说，星系正以由其红移定义的速度“远离我们而去”。在严格的层面上这是无意义的胡话，是 Wittgenstein 所说的“语法错误”——星系并没有在退行，因为“它们相对于我们的速度”这个概念并非良好定义。实际发生的事情是：在我们与星系之间，时空度规沿着光子从这里到那里的路径发生了变化（宇宙膨胀了），导致光的波长增加。作为你怎样出错的一个例子：把多普勒公式天真地应用于星系红移，会推出某些星系的退行速度超光速，表面上与相对论矛盾。这个表观佯谬的解决很简单：星系“退行”这个概念本身就不应当按字面理解。

关于我们不能做的事说得够多了；让我们看看能做什么。平行移动被认为是“保持矢量不变”这个概念在弯曲空间中的推广：即当我们把矢量沿一条路径移动时保持它不变；对任意阶的张量亦然。
给定一条曲线 $x^\mu(\lambda)$，在平直空间中，张量 $T^{\mu_1\mu_2\cdots\mu_k}{}_{\nu_1\nu_2\cdots\nu_l}$ 沿这条曲线保持不变的要求，就是它的分量为常数：
\begin{equation*}
\frac{d}{d\lambda}T^{\mu_1\mu_2\cdots\mu_k}{}_{\nu_1\nu_2\cdots\nu_l} = \frac{dx^\mu}{d\lambda}\frac{\partial}{\partial x^\mu}T^{\mu_1\mu_2\cdots\mu_k}{}_{\nu_1\nu_2\cdots\nu_l} = 0.
\end{equation*}
为了使这一表述具有恰当的张量性，我们只需把这个偏导数换成协变导数，并定义\emph{方向协变导数}（directional covariant derivative）为
\begin{equation*}
\frac{D}{d\lambda} = \frac{dx^\mu}{d\lambda}\nabla_\mu.
\tag{3.38}
\end{equation*}
这是一个只在路径上有定义的映射，从 $(k,l)$ 型张量到 $(k,l)$ 型张量。于是我们把张量 $T$ 沿路径 $x^\mu(\lambda)$ 的\emph{平行移动}（parallel transport）定义为：$T$ 沿路径的协变导数为零这一要求：
\begin{equation*}
\left(\frac{D}{d\lambda}T\right)^{\mu_1\mu_2\cdots\mu_k}{}_{\nu_1\nu_2\cdots\nu_l} \equiv \frac{dx^\sigma}{d\lambda}\nabla_\sigma T^{\mu_1\mu_2\cdots\mu_k}{}_{\nu_1\nu_2\cdots\nu_l} = 0.
\tag{3.39}
\end{equation*}
这是一个良好定义的张量方程（因为切矢量 $dx^\mu/d\lambda$ 和协变导数 $\nabla T$ 都是张量），称为\emph{平行移动方程}（equation of parallel transport）。对矢量而言，它取如下形式：
\begin{equation*}
\frac{d}{d\lambda}V^\mu + \Gamma^\mu_{\sigma\rho}\frac{dx^\sigma}{d\lambda}V^\rho = 0.
\tag{3.40}
\end{equation*}
我们可以把平行移动方程看作一个定义初值问题的一阶微分方程：给定路径上某一点处的张量，沿路径的其他各点将存在该张量的一个唯一的延续，使得这个延续求解式 (3.39)。我们说这样一个张量是被平行移动的（parallel-transported）。

平行移动的概念显然依赖于联络，而不同的联络导致不同的答案。如果联络是度规相容的，那么度规总是相对于它被平行移动的：
\begin{equation*}
\frac{D}{d\lambda}g_{\mu\nu} = \frac{dx^\sigma}{d\lambda}\nabla_\sigma g_{\mu\nu} = 0.
\tag{3.41}
\end{equation*}
由此可得，两个被平行移动的矢量的内积保持不变。也就是说，若 $V^\mu$ 和 $W^\nu$ 沿曲线 $x^\mu(\lambda)$ 被平行移动，则有
\begin{align*}
\frac{D}{d\lambda}(g_{\mu\nu}V^\mu W^\nu) = \left(\frac{D}{d\lambda}g_{\mu\nu}\right)V^\mu W^\nu + g_{\mu\nu}\left(\frac{D}{d\lambda}V^\mu\right)W^\nu + g_{\mu\nu}V^\mu\left(\frac{D}{d\lambda}W^\nu\right)\\
= 0.
\tag{3.42}
\end{align*}
这意味着，相对于度规相容联络的平行移动保持矢量的模（norm）、正交性的判断，等等。

有了平行移动的定义，下一个合乎逻辑的步骤就是讨论测地线。测地线是欧几里得空间中直线概念在弯曲空间中的推广。我们都知道直线是什么：它是距离最短的路径
% 本块末句在原书书页106顶部继续（"between two points"），下一块请用 %JOIN% 续接本句。

